%%% FIELD proof-calendar-lag-loading
Because \(\pi_s(X)\geq c_{\pi}>0\) almost surely for every \(s\in\mathcal S\), every ratio in the definition of \(\theta(\Pi)\) is well defined and bounded by \(c_{\pi}^{-1}\).  The conditional means defining \(\mu_t(s,X)\) ensure integrability of every supported potential-outcome coordinate; finite support, finite \(T\), and \(G(\Pi)\succ0\) then ensure that the expectations and finite interchanges below are well defined.  Write \(\mu(s,X)=(\mu_1(s,X),\ldots,\mu_T(s,X))^{\top}\).  Conditional on \(X\), consistency and pathwise conditional mean random assignment give
\[
\begin{aligned}
\mathbb E_P\!\left[\left.\frac{\Pi_W}{\pi_W(X)}A_W(\Pi)Y\right|X\right]
&=\sum_{s\in\mathcal S}\pi_s(X)\frac{\Pi_s}{\pi_s(X)}A_s(\Pi)
  \mathbb E_P[Y\mid W=s,X]\\
&=\sum_{s\in\mathcal S}\Pi_sA_s(\Pi)\mu(s,X).
\end{aligned}
\]
Indeed, on \(\{W=s\}\), consistency gives \(Y_t=Y_t(s)\), and the random-assignment assumption then gives \(\mathbb E_P[Y_t\mid W=s,X]=\mu_t(s,X)\) for every coordinate \(t\).

Let \(m(X)=(m_1(X),\ldots,m_T(X))^{\top}\).  The calendar-memory assumption is equivalently
\[
 \mu(s,X)=m(X)+\sum_{t=1}^{T}\sum_{\ell=0}^{p}
 e_t\tau_{t,\ell}(X)(D_s)_{t,\ell}.
\]
Since \(\Pi\) is a probability law and \(M(\Pi)=\sum_s\Pi_sD_s\),
\[
 \sum_{s\in\mathcal S}\Pi_sR_s(\Pi)
 =J\left\{\sum_{s\in\mathcal S}\Pi_sD_s-M(\Pi)\sum_{s\in\mathcal S}\Pi_s\right\}=0.
\]
Consequently the unrestricted baseline cancels exactly:
\[
 \sum_{s\in\mathcal S}\Pi_sA_s(\Pi)m(X)
 =T^{-1}G(\Pi)^{-1}
   \left\{\sum_{s\in\mathcal S}\Pi_sR_s(\Pi)^{\top}\right\}m(X)=0.
\]
For every effect cell \((t,\ell)\), the definitions of \(A_s\), \(h_{t,\ell}\), and \(L\) give
\[
\begin{aligned}
 \sum_{s\in\mathcal S}\Pi_sA_s(\Pi)e_t(D_s)_{t,\ell}
 &=G(\Pi)^{-1}\left\{T^{-1}\sum_{s\in\mathcal S}
   \Pi_sR_s(\Pi)^{\top}e_t(D_s)_{t,\ell}\right\}\\
 &=G(\Pi)^{-1}h_{t,\ell}(\Pi)
 =L_{\cdot,t,\ell}(\Pi).
\end{aligned}
\]
The inverse exists by \(G(\Pi)\succ0\).  Taking expectation over \(X\), interchanging only finite sums, and using \(\tau_{t,\ell}=\mathbb E_P[\tau_{t,\ell}(X)]\), proves
\[
 \theta(\Pi)=\sum_{t=1}^{T}\sum_{\ell=0}^{p}
 L_{\cdot,t,\ell}(\Pi)\tau_{t,\ell}.
\]
Thus the causal target array is kept distinct from the observed-law functional \(\theta(\Pi)\), and the displayed equality follows from the stated assignment and outcome-mean conditions rather than by definition.

It remains to derive the design identities.  Regard \(s\) as a random path with law \(\Pi\), and put \(\widetilde D_s=D_s-M(\Pi)\).  Since \(J^{\top}=J\) and \(J^2=J\),
\[
 G(\Pi)=T^{-1}\mathbb E_{\Pi}[\widetilde D_s^{\top}J\widetilde D_s].
\]
Because \((D_s)_{t,j}=s_{t-j}\), with out-of-panel coordinates set to zero, its \((j,r)\) entry is
\[
\begin{aligned}
 G_{jr}(\Pi)
 &=T^{-1}\sum_{u=1}^{T}\sum_{v=1}^{T}J_{uv}
   \operatorname{Cov}_{\Pi}(s_{u-j},s_{v-r})\\
 &=T^{-1}\sum_{t=1}^{T}C_{\Pi}(t-j,t-r)
   -T^{-2}\sum_{t=1}^{T}\sum_{v=1}^{T}C_{\Pi}(t-j,v-r),
\end{aligned}
\]
where \(J_{uv}=\mathbf 1\{u=v\}-T^{-1}\).  Similarly, centering the first factor and observing that its product with the uncentered second factor has the same expectation as the corresponding covariance,
\[
\begin{aligned}
 h_{t,\ell,j}(\Pi)
 &=T^{-1}\sum_{v=1}^{T}J_{vt}
   \operatorname{Cov}_{\Pi}(s_{v-j},s_{t-\ell})\\
 &=T^{-1}\left\{C_{\Pi}(t-j,t-\ell)
   -T^{-1}\sum_{v=1}^{T}C_{\Pi}(v-j,t-\ell)\right\}.
\end{aligned}
\]
These two calculations require only the finite binary-path support.

If every supported path is absorbing, then for every coordinate occurring above,
\(s_u=\mathbf 1\{a(s)\leq u\}\), with value zero for \(u\leq0\).  Hence
\[
\begin{aligned}
 C_{\Pi}(u,v)
 &=\mathbb E_{\Pi}[s_us_v]-\mathbb E_{\Pi}[s_u]\mathbb E_{\Pi}[s_v]\\
 &=F_{\Pi}(\min\{u,v\})-F_{\Pi}(u)F_{\Pi}(v)
 =K_{\Pi}(u,v).
\end{aligned}
\]
This proves the ordered-Green specialization.

Finally, let \(e_{\ell}\in\mathbb R^d\).  Summing the numerator over calendar time gives
\[
\begin{aligned}
 \sum_{t=1}^{T}h_{t,\ell}(\Pi)
 &=T^{-1}\sum_{s\in\mathcal S}\Pi_sR_s(\Pi)^{\top}D_se_{\ell}\\
 &=T^{-1}\sum_{s\in\mathcal S}\Pi_sR_s(\Pi)^{\top}
   \{D_s-M(\Pi)\}e_{\ell}\\
 &=T^{-1}\sum_{s\in\mathcal S}\Pi_sR_s(\Pi)^{\top}R_s(\Pi)e_{\ell}\\
 &=G(\Pi)e_{\ell}.
\end{aligned}
\]
The second equality uses \(\sum_s\Pi_sR_s(\Pi)=0\); the third uses \(R_s=J\widetilde D_s\), \(J^{\top}=J\), and \(J^2=J\).  Multiplication by \(G(\Pi)^{-1}\) therefore yields
\[
 \sum_{t=1}^{T}L_{j,t,\ell}(\Pi)
 =e_j^{\top}e_{\ell}=\mathbf 1\{j=\ell\},
\]
as claimed.
%%% FIELD proof-homogeneous-strict-extension
Assume calendar homogeneity and denote the common effect at lag \(\ell\) by \(\overline\tau_{\ell}\), so that \(\tau_{t,\ell}=\overline\tau_{\ell}\) for every legal calendar index.  The causal loading representation in \(\mathrm{thm{:}calendar\text{-}lag\text{-}loading}\), followed by its normalization identity, gives for every \(j\in\{0,\ldots,p\}\)
\[
\begin{aligned}
 \theta_j(\Pi)
 &=\sum_{\ell=0}^{p}\sum_{t=1}^{T}
   L_{j,t,\ell}(\Pi)\overline\tau_{\ell}\\
 &=\sum_{\ell=0}^{p}\mathbf 1\{j=\ell\}\overline\tau_{\ell}
 =\overline\tau_j.
\end{aligned}
\]
Thus \(\theta_j(\Pi)=\tau_{t,j}\) for every legal \(t\).  This conclusion is coefficient algebra in the present superpopulation model.  Independently, \(\mathrm{thm{:}published\text{-}class\text{-}calendar\text{-}loading}\) applies under the published unit-indexed known-propensity and weak-dependence conditions and, under the same calendar homogeneity restriction, gives the published-class recovery statement with remainder \(O_P(n^{-q/2})\).  Therefore no class-inclusion argument is used.

For the strict calendar-varying witness, take \(T=3\), \(p=1\), \(\mathcal S=\{000,001,011\}\), and put mass \(1/3\) on each listed path.  In that listing order,
\[
 D_{000}=\begin{pmatrix}0&0\\0&0\\0&0\end{pmatrix},\qquad
 D_{001}=\begin{pmatrix}0&0\\0&0\\1&0\end{pmatrix},\qquad
 D_{011}=\begin{pmatrix}0&0\\1&0\\1&1\end{pmatrix},
\]
and hence
\[
 M(\Pi)=\begin{pmatrix}0&0\\1/3&0\\2/3&1/3\end{pmatrix},\qquad
 J=\begin{pmatrix}2/3&-1/3&-1/3\\-1/3&2/3&-1/3\\-1/3&-1/3&2/3\end{pmatrix}.
\]
Substitution into \(R_s=J(D_s-M(\Pi))\), followed by finite matrix multiplication in the definitions of \(G\) and \(h\), gives
\[
 G(\Pi)=\begin{pmatrix}2/27&0\\0&4/81\end{pmatrix},
\]
\[
 h_{2,0}(\Pi)=\begin{pmatrix}1/27\\-2/81\end{pmatrix},\qquad
 h_{3,0}(\Pi)=\begin{pmatrix}1/27\\2/81\end{pmatrix},\qquad
 h_{3,1}(\Pi)=\begin{pmatrix}0\\4/81\end{pmatrix},
\]
with every other \(h_{t,\ell}(\Pi)\) equal to zero.  In particular, \(G(\Pi)\succ0\), and multiplication by
\(G(\Pi)^{-1}=\operatorname{diag}(27/2,81/4)\) yields
\[
 L_{\cdot,2,0}(\Pi)=\begin{pmatrix}1/2\\-1/2\end{pmatrix},\qquad
 L_{\cdot,3,0}(\Pi)=\begin{pmatrix}1/2\\1/2\end{pmatrix},\qquad
 L_{\cdot,3,1}(\Pi)=\begin{pmatrix}0\\1\end{pmatrix}.
\]
The already-proved causal loading representation now gives
\[
 \theta_0=\frac{\tau_{2,0}+\tau_{3,0}}{2},\qquad
 \theta_1=\frac{-\tau_{2,0}+\tau_{3,0}}{2}+\tau_{3,1}.
\]
Thus, holding \(\tau_{3,1}\) fixed while varying either calendar-specific current effect changes the lag-one coefficient, so no time-homogeneous lag-effect vector determines the coefficient on this calendar-varying class.

The same witness has an independent randomized-assignment embedding in the published sampling setup.  For each unit, draw \(W_i\) independently from \(\Pi\), and define bounded deterministic potential outcomes for every binary path \(w\) by
\[
 Y_{it}(w)=\sum_{\ell=0}^{1}\tau_{t,\ell}w_{t-\ell},
 \qquad w_u=0\text{ outside }\{1,2,3\}.
\]
Then the propensity conditional on the full potential-outcome schedule equals \(\pi_i(s)=1/3\) for \(s\in\mathcal S\), so published overlap holds with \(c_*=1/3\).  Independence gives \(\rho_{ij}=0\) for \(i\neq j\), while \(\rho_{ii}=1\), and therefore
\[
 n^{-2}\sum_{i=1}^{n}\sum_{j=1}^{n}\rho_{ij}=n^{-1}.
\]
The finite collection of deterministic outcomes has a uniform finite second-moment bound, and the displayed \(G(\Pi)\) supplies full rank.  Hence all published regularity conditions hold with \(q=1\), and \(\mathrm{thm{:}published\text{-}class\text{-}calendar\text{-}loading}\) transfers the same calendar-varying nonreduction to the published estimator up to its \(O_P(n^{-1/2})\) remainder.  Under calendar homogeneity, that same theorem instead gives the exact published Theorem \(4.1\) recovery conclusion.
%%% FIELD proof-adjacent-ratio-determinant-sign-rule
Because an absorbing binary path is determined by its adoption date, distinct supported paths have distinct adoption dates.  Under the stated ordering, if \(u_a=a(s_a)<u_b=a(s_b)\), then the boundary convention implies for every row \(r\) and lag column \(j\)
\[
\begin{aligned}
 (E_{ab})_{r,j}
 &=(D_{s_a})_{r,j}-(D_{s_b})_{r,j}\\
 &=\mathbf 1\{u_a\leq r-j\}-\mathbf 1\{u_b\leq r-j\}\\
 &=\mathbf 1\{u_a\leq r-j<u_b\}
 =(I_{ab})_{r,j}.
\end{aligned}
\]
Thus \(E_{ab}=I_{ab}\), including coordinates shifted outside the panel.  The established absorbing graph-rank lemma gives
\[
 \ker H_{\mathcal S}=\{0\}
 \quad\Longleftrightarrow\quad
 \text{every lag vertex of }\mathfrak G_{\mathcal S,p}\text{ is connected to }\bot.
\]

The floor \(\Pi_{s_a}\geq\epsilon>0\) makes every adjacent ratio \(\rho_a\) strictly positive and well defined.  Recursion in the ratios and normalization of \(\Pi\) give
\[
 \frac{\Pi_{s_a}}{\Pi_{s_1}}
 =\prod_{r=1}^{a-1}\rho_r=z_a,
 \qquad
 1=\sum_{a=1}^{k}\Pi_{s_a}=\Pi_{s_1}Z,
 \qquad
 \Pi_{s_a}=\frac{z_a}{Z}.
\]
Using \(E_{ab}=I_{ab}\) in the pair contributions and then applying the established pair-mass identities yields
\[
\begin{aligned}
 G(\Pi)
 &=\sum_{a<b}\Pi_{s_a}\Pi_{s_b}B_{ab}
 =Z^{-2}\Gamma(\rho),\\
 h_{t,\ell}(\Pi)
 &=\sum_{a<b}\Pi_{s_a}\Pi_{s_b}b_{ab,t,\ell}
 =Z^{-2}g_{t,\ell}(\rho).
\end{aligned}
\]
By construction \(G(\Pi)\) is symmetric positive semidefinite.  Positivity of all masses and the support-rank lemma give
\(\ker G(\Pi)=\ker H_{\mathcal S}=\{0\}\); hence \(G(\Pi)\succ0\).  It follows that
\[
 \Gamma(\rho)=Z^2G(\Pi)\succ0,
 \qquad \det\Gamma(\rho)>0,
\]
and the common positive scale cancels from the loading:
\[
 L_{\cdot,t,\ell}(\Pi)=\Gamma(\rho)^{-1}g_{t,\ell}(\rho).
\]

For completeness, the required Cramer identity follows from determinant expansion.  Let \(\operatorname{adj}(\Gamma)\) be the cofactor adjugate.  Cofactor expansion gives
\[
 \Gamma\operatorname{adj}(\Gamma)=\det(\Gamma)I_d:
\]
on a diagonal entry this is the determinant expansion, whereas an off-diagonal entry is the same expansion for a matrix with two equal columns and is zero.  Expanding the determinant of the matrix formed by replacing column \(j\) of \(\Gamma\) by \(g_{t,\ell}\) gives
\[
 e_j^{\top}\operatorname{adj}(\Gamma)g_{t,\ell}
 =\Delta_{j,t,\ell}.
\]
Since \(\det\Gamma>0\), these identities imply
\[
 L_{j,t,\ell}(\Pi)
 =\frac{e_j^{\top}\operatorname{adj}(\Gamma(\rho))g_{t,\ell}(\rho)}
        {\det\Gamma(\rho)}
 =\frac{\Delta_{j,t,\ell}(\rho)}{\det\Gamma(\rho)}.
\]
The denominator is strictly positive, so
\[
 \operatorname{sgn}L_{j,t,\ell}(\Pi)
 =\operatorname{sgn}\Delta_{j,t,\ell}(\rho),
 \qquad
 L_{j,t,\ell}(\Pi)=0
 \Longleftrightarrow
 \Delta_{j,t,\ell}(\rho)=0.
\]
Taking the finite conjunction over all \(t\) and all \(j\neq\ell\) proves the stated characterization of vanishing off-diagonal causal-lag loadings.

Finally, \(\Pi_{s_a}=z_a/Z\) proves, in both directions,
\[
 \Pi_{s_a}\geq\epsilon\text{ for every }a
 \quad\Longleftrightarrow\quad
 z_a\geq\epsilon Z\text{ for every }a.
\]
Each \(z_a\) is a finite monomial in the adjacent ratios, while every endpoint \(u_a\) is an integer and every entry of \(I_{ab}\) and \(J\) is rational.  Thus, on rationally represented ratios, all entries of \(\Gamma\) and \(g\), the floor comparisons, and all determinants \(\Delta_{j,t,\ell}\) are computable by finitely many exact rational operations.  Exact determinant evaluation and comparison with zero terminate, including for shifted boundary coordinates.  This supplies a pointwise certificate for the given law; the argument does not enumerate the global semialgebraic sign chambers over the admissible ratio domain.
